\documentclass[10.5pt,letterpaper]{article}
\usepackage[margin=0.58in,top=0.52in,bottom=0.52in]{geometry}
\usepackage[T1]{fontenc}
\usepackage{newtxtext}
\usepackage{xcolor}
\usepackage{enumitem}
\usepackage{hyperref}
\definecolor{sectionblue}{HTML}{25486E}
\hypersetup{colorlinks=true,urlcolor=sectionblue}
\pagestyle{plain}
\setlength{\parindent}{0pt}
\setlist[itemize]{leftmargin=1.1em,itemsep=1.4pt,topsep=2pt,parsep=0pt}
\newcommand{\ressection}[1]{\vspace{6pt}{\Large\bfseries\color{sectionblue}#1}\par\vspace{-2pt}\hrule\vspace{4pt}}
\begin{document}
{\Huge\bfseries Noah Csaky-Schwede}\par
{\large\bfseries Master Resume - Source Document for Tailored Applications}\par
Mississauga, ON / Victoria, BC \quad | \quad 650-284-9838 \quad | \quad noah.csakyschwede@mail.utoronto.ca \quad | \quad github.com/ncsaky\par
\vspace{6pt}
\textit{Note: This master resume is intentionally comprehensive and should be trimmed, reordered, and rewritten for each specific GIS, CS, software, data, or AI-focused job posting.}
\ressection{Target Profile}
Honours BSc student at the University of Toronto Mississauga majoring in Geographical Information Systems with minors in Computer Science and Mathematical Sciences. Interdisciplinary background across GIS, spatial analysis, software engineering, systems programming, data science, machine learning, cartographic design, and technical education. Particularly interested in the future of software development through AI-assisted and agentic coding workflows, including human-in-the-loop review, iterative prompting, code generation, debugging, testing, and integration.
\ressection{AI-Assisted and Agentic Software Development Focus}
\begin{itemize}
\item Deep practical interest in agentic software development and LLM-assisted coding workflows, with years of using ChatGPT as a programming, debugging, design, and learning tool throughout university.
\item Hands-on experience using Codex/Codex CLI, Claude Code CLI and app workflows, Qwen Coder CLI, local Gemma-family model experiments, and Hermes-style agent harnesses for software development exploration.
\item Applied AI-assisted development with human review to software coursework, independent projects, debugging, codebase exploration, implementation planning, refactoring, and documentation workflows.
\item Interested in the shift from manual coding toward agent-directed software production, where developers define requirements, validate outputs, review architecture, test behavior, and integrate generated code responsibly.
\item Familiar with the broader ecosystem of emerging coding agents and open-source agentic tools, including CLI-first agent workflows and local/self-hosted model experimentation.
\item Strong emphasis on treating AI-generated code as draft material requiring verification, testing, security awareness, architectural review, and human ownership.
\end{itemize}
\ressection{Education}
\textbf{University of Toronto Mississauga}\hfill Mississauga, ON\par
\textit{Honours Bachelor of Science}\hfill Expected completion: Final semester remaining\par
Major: Geographical Information Systems | Minors: Computer Science, Mathematical Sciences
\ressection{Skills Bank}
\begin{tabular}{p{1.65in}p{5.35in}}
\textbf{Programming} & Python, Java, C, JavaScript, TypeScript, R, SQL, HTML, CSS. \\[-1pt]
\textbf{Software Engineering} & React, Next.js, Flask, REST APIs, Git/GitHub, object-oriented programming, software design patterns, agile development, SDLC, requirements planning, unit testing, debugging, refactoring, code review, API integration. \\[-1pt]
\textbf{AI/Agentic Development} & ChatGPT-assisted coding, Codex/Codex CLI, Claude Code CLI/app workflows, Qwen Coder CLI exploration, local Gemma-family model experimentation, agent harness exploration, prompt-driven implementation, human-in-the-loop review, test-and-verify workflows. \\[-1pt]
\textbf{GIS and Spatial Analysis} & QGIS, ArcGIS Pro, digital cartography, thematic mapping, coordinate systems, map projections, spatial joins, buffer/proximity analysis, vector data, raster data, geospatial visualization, environmental spatial analysis. \\[-1pt]
\textbf{Remote Sensing} & Satellite imagery, raster layers, image interpretation, environmental remote sensing, sensor systems, land-surface analysis, multi-layer imagery workflows, vegetation/water/soil interpretation. \\[-1pt]
\textbf{Data Science and ML} & pandas, scikit-learn, TPOT AutoML, regression, classification, clustering, model evaluation, feature engineering, exploratory data analysis, Selenium scraping, JSON processing, multiprocessing, ensemble models. \\[-1pt]
\textbf{Systems and Tools} & Linux/Unix, macOS, Windows, VS Code, Excel, command-line tools, shell utilities, file processing, process management, pipes, signals, C memory management. \\[-1pt]
\end{tabular}
\ressection{Experience}
\textbf{Program Facilitator and Content Developer} \hfill \textit{University of Toronto Engineering Outreach - UofT Co-Op Program / Summer 2024 / Toronto, ON}\par
\begin{itemize}
\item Designed and delivered a week-long app design course for high school students using Python, Flask, HTML, and CSS.
\item Developed curriculum introducing web development, application design, debugging, and project-based programming.
\item Iteratively refined course materials across multiple sessions based on student feedback and observed learning gaps.
\item Guided students through building functional web applications while encouraging problem-solving, creativity, and independent debugging.
\item Communicated technical concepts to beginner programmers in an accessible and structured way.
\item Balanced classroom management, technical support, and curriculum delivery in a fast-paced educational environment.
\end{itemize}
\textbf{Educational Software Lead} \hfill \textit{Byte Camp / Summer 2023 / Richmond, BC}\par
\begin{itemize}
\item Taught programming fundamentals including loops, conditionals, functions, and object-oriented programming to young learners.
\item Designed engaging lesson plans and coding projects that balanced technical depth with creativity and collaboration.
\item Supported students through debugging, project planning, and iterative development.
\item Created an interactive learning environment focused on teamwork, problem-solving, and confidence with technology.
\item Developed communication and leadership skills through daily instruction and technical mentoring.
\end{itemize}
\ressection{Major Software, Data, GIS, and AI Projects}
\textbf{UTM Campus Dashboard} \hfill \textit{Semester-Long Software Engineering Project / Full-Stack Development / AI-Assisted Development}\par
\begin{itemize}
\item Built a full-stack campus services platform over a semester-long CSC301 software engineering project as part of a seven-person team using agile sprints, recurring scrums, Jira issue tracking, Slack coordination, Git branches, pull requests, and teammate code review.
\item Helped develop a centralized UTM campus dashboard for students to check food wait times, gym crowding, parking availability, campus events, lost-and-found posts, and building-level campus information.
\item Contributed to a Next.js/React/TypeScript frontend with reusable components, route-based pages, modals, cards, charts, dark-mode/accessibility UI improvements, and interactive user workflows.
\item Worked across a Flask backend using SQLAlchemy, SQLite, REST-style API routes, Flask-CORS, Flask-Bcrypt, session-based authentication, and database-backed models for users, reports, events, lost-and-found posts, and planner data.
\item Implemented and debugged frontend-backend report submission flows so logged-in users could submit campus status reports from food, gym, and parking interfaces, with backend routes receiving and storing report data.
\item Built an interactive UTM campus map feature with clickable building regions, color-coded occupancy/status indicators, building detail panels, and navigation links into related food, gym, parking, and events pages.
\item Developed a personalized day-planner feature that parses uploaded .ics course calendar files and generates daily schedules around classes, meals, gym time, study blocks, commute timing, campus events, and user preferences.
\item Added planner interactions including proportional event blocks, past-event handling, custom user events, editable schedule items, saved planner state, conflict handling, and manual adjustments for demo-ready usability.
\item Supported feature development beyond the initial campus-status concept, including event calendar views, event shuffle interactions, lost-and-found CRUD workflows, Bluetooth/Wi-Fi scanner integration experiments, and seeded demo data for presentations.
\item Used AI-assisted and agentic coding workflows as part of the development process, including codebase exploration, implementation planning, debugging, branch-diff review, documentation drafting, local testing guidance, and iterative feature refinement.
\item Practiced responsible human-in-the-loop AI usage by reviewing generated code, testing behavior locally, resolving merge conflicts, constraining changes to ticket scope, and coordinating PR visibility with teammates before integration.
\item Produced sprint deliverables including sprint retrospectives, user-story quality evaluations, technical handoff documentation, PR summaries, demo notes, and written explanations of architecture and implementation decisions.
\item Technologies: Next.js, React, TypeScript, Tailwind CSS, Recharts, Flask, SQLAlchemy, SQLite, Python, REST APIs, Flask-CORS, Flask-Bcrypt, Git, GitHub, Jira, Slack, AI-assisted development.
\end{itemize}
\textbf{Agentic Basketball Simulation and Analytics Platform} \hfill \textit{Independent Project / Agentic Software Development / Sports Simulation}\par
\begin{itemize}
\item Building a CLI-first NBA general-manager simulation platform inspired by and intended to improve on franchise-management modes such as NBA 2K GM Mode, combining sports analytics, stochastic simulation, roster-building AI, financial decision-making, scouting, staff management, narrative systems, and multi-season league saves.
\item Directed development through long-running human-in-the-loop Codex sessions: wrote detailed product specifications and milestone plans, decomposed a large ambiguous game vision into implementable subsystems, reviewed generated code, diagnosed playtest failures, constrained architecture, ran regression tests, and iterated on simulation realism and user experience.
\item Grew the project into a modular Python codebase with 21 package modules and 42,128 physical lines across source and tests, including separate layers for schema/data modelling, ingestion, auditing, canonical storage, simulation, save-state mutation, health/development, transactions, contracts, draft/scouting, staff management, narrative generation, loading-screen animation, research data, and the interactive play loop.
\item Designed a canonical 2025-26 NBA preseason data foundation using deterministic JSON and SQLite exports, with records for players, teams, traits, contracts, draft picks, roster slots, staff profiles, source evidence, confidence scores, and coverage reports.
\item Built ingestion, override, and calibration workflows that treat raw files and public research as inputs rather than truth, allowing manual corrections, confidence notes, source citations, deterministic rebuilds, subjective full-health rating priors, display-scale ratings, and separation between immutable preseason canonical data and mutable save-state consequences.
\item Modelled hidden basketball traits as first-class data, including shooting range, release speed, shot versatility, handle pressure, rim pressure, passing reads, offensive rebounding, defensive effort, lateral agility, screen navigation, rim deterrence, scheme IQ, stamina, portability, playoff translation, and health risk.
\item Implemented an availability-aware simulation engine with two distinct modes: real-minutes replay for validation against actual 2025-26 NBA games and sandbox simulation for alternate-universe saves, fantasy-style rosters, injuries, trades, development, playoffs, draft classes, and future seasons.
\item Added possession and player-output modelling for scoring, assists, rebounds, steals, blocks, turnovers, shooting lines, free throws, three-point volume, health/fatigue/rust effects, overtime outcomes, team context, coaching modifiers, player-line variance, and nonlinear lineup interactions such as elite spacing, rim deterrence, defensive weak links, creation burden, and offensive rebounding.
\item Developed market-aware calibration tooling using historical betting inputs where available, Monte Carlo validation, no-vig win probabilities, Brier score, log loss, spread/total comparison, real-outcome percentile checks, standings sanity checks, awards calibration, and market-versus-simulation disagreement reports.
\item Built a persistent league\_save\_v1 format with user team, current date, league phase, seed, rosters, contracts, picks, staff slots, health states, injuries, player development, transaction logs, standings, game results, box scores, season totals, news items, social feed items, narrative cache, starting lineups, and reviewable AI trade offers.
\item Implemented a playable calendar and phase engine with preseason, regular season, trade deadline, play-in/playoffs, draft lottery, draft, free agency, training camp, legal-action gates, scheduled game simulation, standings, league leaders, team dashboards, box-score viewing, and multi-season rollover.
\item Built explainable AI-GM systems for trades, extensions, free agency, released-player signings, draft decisions, draft-night trades, free-agency trades, and staff changes using team strategic state, front-office personality, competence, pressure, timeline, roster needs, cap posture, youth pipeline, pick inventory, player age curves, contract surplus/drag, health risk, role scarcity, and playoff fit.
\item Implemented practical NBA-style transaction logic, including salary matching, roster-count effects, recently traded/signed restrictions, unresolved-salary manual review, pick ownership, current-draft pick identity, conditional firsts, protection fallback collateral, pick swaps, used-pick restrictions, and Stepien-style first-round guardrails while intentionally deferring deep CBA edge cases.
\item Created trade-finder and trade-builder workflows that generate legal multi-asset packages, support user-controlled trade negotiation, prompt for protected-pick or pick-swap terms before valuation, expose exact legality failures, and display package-value breakdowns across player quality, age/timeline, contract value, lineup fit, health, picks, cap flexibility, and GM personality/pressure modifiers.
\item Built contract-market and negotiation systems for extensions and free agency, estimating salary asks/offers from league salary comps, player value, role tier, age, health, upside, playoff value, team context, cap posture, GM personality, and player preferences for money, role, winning, security, loyalty, and fit; added AI extension timing, refusals, trade-demand pressure, and staff/player retention windows.
\item Implemented draft systems including real 2026 prospect ingestion, generated future draft classes, scouting fog, scouting-staff influence, prospect boards, pick recommendations, rookie contract projection, lottery/order handling, draft-night trade markets, rookie rotation onboarding, and AI draft logic balancing best-player-available with team need, fit, upside, cap value, redundancy penalties, and GM risk tolerance.
\item Added staff management with fictional gameplay staff slots for head coach, offensive coordinator, defensive coordinator, development lead, scouting lead, and performance lead; staff contracts, budgets, negotiations, firings, replacements, interim staff, market candidates, AI staff-retention logic, role-specific firing triggers, public league events, and system effects on simulation, development, scouting, health, and fatigue.
\item Built lightweight sandbox injury, fatigue, recovery, rust, and monthly player-development systems using player health profiles, age/workload/stamina risk, generated injury events, performance-staff modifiers, trait-level growth/regression, startup injury handling, in-season free-agent signing repair, and annual development review displays.
\item Designed CLI/TUI-style user experience features including numbered interactive play sessions, save picker, random team creation, loading-screen video-to-ASCII animation, team dashboards, ratings guide, starting-five visualization, standings, calendar/box-score views, draft room, free-agency room, staff room, trade room, playoff dashboard, league-events browser, morale bars, and readable transaction/news recaps.
\item Added an optional local Ollama-backed narrative engine that lazily rewrites social posts and generates press-conference mini-games from sandbox-only context packets, with persona styles, JSON validation, deterministic caching, fallback templates, factual guardrails, cap/health/rating evidence snippets, and no dependency on outside NBA memory.
\item Used repeated playtest feedback loops to identify and fix complex behavioural bugs such as illegal trades, rejected trade-finder offers, unrealistic trade packages, duplicated news, stale free-agent pools, protected-pick fallback errors, lottery/draft-order mismatches, roster-cut ghost stats, staff-market churn, player-rating anomalies, unclear dashboards, slow hot paths, and phase-inappropriate menu actions.
\item Maintained a regression-focused workflow with Python unittest coverage; current local suite contains 166 passing tests covering deterministic canonical builds, schema validation, health/development, simulation, market calibration, trade AI, contract/free-agency logic, draft systems, pick obligations/swaps, staff management, save advancement, social/press narrative behaviour, and multi-season rollover.
\item Technologies and practices: Python 3.11, dataclasses, JSON, SQLite, CLI design, deterministic seeded simulation, data ingestion, sports analytics, Monte Carlo validation, optional local LLM integration, terminal animation, testing, debugging, product specification, agent-directed development, human-in-the-loop code review, and iterative UX/playtest refinement.
\end{itemize}
\textbf{Environmental Equity and Roadway Exposure Research Proposal} \hfill \textit{GIS / Urban Geography / Environmental Justice}\par
\begin{itemize}
\item Developed a GIS-focused research proposal examining whether selected demographic groups in Mississauga are disproportionately located near major roadways.
\item Connected spatial inequality, commuter-oriented urban form, and emerging research on airborne microplastics and tire-wear particles as potential exposure risks near high-traffic corridors.
\item Proposed a spatial workflow using roadway proximity buffers, demographic overlays, census-based variables, and map-based interpretation to evaluate environmental equity concerns.
\item Framed the project around public-health geography, environmental justice, and the uneven distribution of urban infrastructure impacts.
\item Designed the proposal to investigate whether transportation infrastructure and residential distribution patterns create unequal exposure risks across population groups.
\item Relevant skills: GIS, environmental justice, urban geography, spatial analysis, demographic analysis, proximity analysis, research design.
\end{itemize}
\textbf{Environmental Remote Sensing and Raster Analysis} \hfill \textit{GIS / Satellite Imagery / Environmental Analysis}\par
\begin{itemize}
\item Worked with remotely sensed imagery and multi-layer raster datasets to analyze environmental and land-surface patterns.
\item Applied remote sensing concepts including sensor systems, image interpretation, surface-energy interactions, and environmental monitoring.
\item Interpreted satellite imagery for applications related to vegetation, water, soil, and broader land-resource analysis.
\item Built experience managing spatial data layers, comparing imagery, and translating raster-based observations into geographic conclusions.
\item Strengthened ability to work with image-based spatial data and connect visual patterns to environmental processes.
\item Relevant skills: remote sensing, raster data, satellite imagery, environmental analysis, image interpretation, GIS data layers.
\end{itemize}
\textbf{Digital Mapping and Cartographic Design Portfolio} \hfill \textit{GIS / Cartography / Thematic Mapping}\par
\begin{itemize}
\item Produced GIS-based maps using cartographic principles, coordinate systems, map projections, classification methods, and spatial measurement techniques.
\item Designed thematic maps with attention to visual hierarchy, legend structure, scale, projection choice, symbology, and map readability.
\item Applied GIS workflows to prepare spatial data, simplify features, organize layers, and communicate mapped patterns clearly.
\item Critically evaluated mapped information and considered how design choices influence interpretation.
\item Built practical experience with map layout, geospatial data preparation, and publication-quality cartographic design.
\item Relevant skills: GIS, QGIS, cartography, thematic mapping, projections, classification, map design, spatial visualization.
\end{itemize}
\textbf{Chess Accuracy Prediction via Stockfish and AutoML} \hfill \textit{Machine Learning / Data Engineering}\par
\begin{itemize}
\item Built an end-to-end machine learning pipeline to predict Chess.com's proprietary Accuracy rating using public gameplay data, PGN files, and Stockfish engine evaluations.
\item Collected large-scale gameplay datasets from authenticated and public Chess.com accounts using Selenium, scraping workflows, and API parsing.
\item Engineered features from PGN data using Stockfish to classify move quality, calculate centipawn loss, and identify game-specific heuristics.
\item Conducted exploratory data analysis to evaluate feature distributions, correlations, and predictive value.
\item Applied TPOT AutoML to train, optimize, and export regression pipelines, including stacked models using ElasticNet, LassoLarsCV, PolynomialFeatures, and RandomForest.
\item Tested feature combinations using lightweight TPOT runs, evaluated mean squared error, R-squared, and adjusted R-squared, then retrained top-performing feature sets with deeper TPOT configurations.
\item Designed and tested an experimental stacked supermodel combining top TPOT-exported pipelines through StackingRegressor.
\item Managed a modular codebase across scraping, feature engineering, model training, evaluation, multiprocessing, and documentation.
\item Technologies: Python, TPOT, scikit-learn, Stockfish, Selenium, pandas, Excel, AutoML, multiprocessing, regression metrics.
\end{itemize}
\textbf{Custom Unix Shell} \hfill \textit{Systems Programming / C / Unix}\par
\begin{itemize}
\item Developed a functional Unix shell supporting job control, foreground and background process execution, I/O redirection, and command pipelines.
\item Implemented signal handling for SIGINT, SIGTSTP, and SIGCHLD to manage process behavior and job state.
\item Designed built-in commands including jobs, bg, fg, and quit.
\item Added support for parsing and executing complex command pipelines.
\item Demonstrated understanding of Unix/Linux internals, process management, system calls, and systems programming in C.
\item Technologies: C, Linux/Unix, process control, signals, pipes, file descriptors.
\end{itemize}
\textbf{Audio Processing Tools in C} \hfill \textit{Systems Programming / Binary File Processing}\par
\begin{itemize}
\item Built command-line tools for processing WAV audio files through binary file manipulation.
\item Implemented remvocals to remove vocals from stereo WAV files by processing channel data.
\item Implemented addecho to add echo effects to mono WAV files using dynamic memory allocation and circular buffering.
\item Used getopt for argument parsing, efficient file I/O, and strict WAV file format handling.
\item Technologies: C, Linux/Unix, binary file processing, memory management, command-line tools.
\end{itemize}
\textbf{NBA Play-by-Play Scraper and Matchup Parser} \hfill \textit{Sports Analytics / Web Scraping / Data Processing}\par
\begin{itemize}
\item Created a Python-based system to scrape NBA play-by-play data from ESPN and identify player matchups across game events.
\item Built a Selenium scraper to extract multi-quarter play-by-play data and organize event outputs into structured JSON.
\item Applied parsing logic to compute on-court durations for specific player lineups and matchup combinations.
\item Strengthened experience with web scraping, sports analytics, event data, and structured data processing.
\item Technologies: Python, Selenium, JSON, data parsing, sports analytics.
\end{itemize}
\ressection{Coursework Bank - Low Priority, Use Selectively}
This section is intentionally a bank for tailoring. For real applications, include only the most relevant courses or remove entirely if projects and skills are stronger.
\textbf{Computer Science}\par
\begin{itemize}
\item CSC148 - Introduction to Computer Science: abstract data types, recursion, program efficiency, object-oriented programming, Python.
\item CSC207 - Software Design: agile methodology, version control, unit testing, design patterns, Java programming, regular expressions, advanced IDE usage.
\item CSC209 - Software Tools and Systems Programming: Unix/Linux development, C programming, shell utilities, processes, memory model, system calls, file processing, pipes, and signals.
\item CSC236 - Introduction to the Theory of Computation: correctness proofs, recursive and iterative algorithms, Master Theorem, automata, formal languages.
\item CSC263 - Data Structures and Analysis: algorithm analysis, abstract data types, data structures, implementation tradeoffs, worst-case and average-case analysis.
\item CSC301 - Introduction to Software Engineering: software engineering methods, agile development, requirements, team-based development, architecture, testing, refactoring, and large group project work.
\item CSC311 - Introduction to Machine Learning: classification, regression, clustering, neural networks, overfitting, generalization, mathematical foundations, and model evaluation.
\end{itemize}
\textbf{GIS, Geography, Remote Sensing, and Spatial Analysis}\par
\begin{itemize}
\item GIS major coursework covering spatial data management, geographic analysis, map production, environmental applications, and geographic research design.
\item Environmental Remote Sensing: satellite imagery, remote sensing systems, image interpretation, environmental monitoring, and raster-based analysis.
\item Digital Mapping and GIS: cartography, thematic mapping, coordinate systems, projections, classification, map production, and critical evaluation of spatial information.
\item Advanced GIS / Research Proposal: spatial research proposal examining environmental equity, roadway exposure, microplastics/tire-wear particles, and demographic distribution in Mississauga.
\end{itemize}
\textbf{Mathematics and Statistics}\par
\begin{itemize}
\item Mathematical Sciences minor coursework involving algorithmic reasoning, applied mathematical thinking, statistical analysis, and quantitative problem-solving.
\item Relevant quantitative background for machine learning, spatial analysis, regression, data modelling, and algorithmic problem solving.
\end{itemize}
\textbf{Leadership, Communication, and Extracurricular}\par
\begin{itemize}
\item Lifelong hockey player with experience in competitive team environments requiring discipline, communication, resilience, and role-specific responsibility.
\item Experience teaching and mentoring younger students in technical subjects through programming camps and university outreach programs.
\item Strong background balancing technical coursework, independent projects, team software development, and applied GIS research.
\item Comfortable explaining technical ideas to non-expert audiences through teaching, project walkthroughs, and collaborative development settings.
\end{itemize}
\ressection{Tailoring Notes - Remove Before Applying}
\begin{itemize}
\item For GIS Analyst roles: prioritize GIS/spatial analysis skills, remote sensing, environmental equity project, cartographic portfolio, Python/R/data analysis, and map production. Cut most low-level systems programming unless the posting asks for programming depth.
\item For Software Developer roles: prioritize Campus Dashboard, Chess ML pipeline, Unix shell, NBA scraper, agentic development focus, Flask/React/Next.js, Python/Java/C/TypeScript, Git, and testing. Keep GIS as a domain specialization.
\item For Data Analyst or ML roles: prioritize Chess Accuracy Prediction, CSC311/ML background, pandas/scikit-learn/TPOT, GIS environmental analysis, remote sensing, scraping, and structured data workflows.
\item For AI/agentic tooling roles: emphasize Codex/Codex CLI, Claude Code, ChatGPT-assisted coding, Qwen Coder CLI, local model experimentation, human-in-the-loop review, testing discipline, and the basketball simulation project.
\end{itemize}
\end{document}
